% 主论文(中文版):偏好坍缩可利用度
% 目标会议:NeurIPS / AAAI 2026
% 编译方式:xelatex main_zh.tex -> bibtex main_zh -> xelatex -> xelatex
\documentclass{article}

% 使用 NeurIPS 2024 样式(按需调整)
\usepackage[final]{neurips_2024}

% 中文支持(XeLaTeX + ctex;使用 TeX Live 自带字体以避免依赖本机 macOS 字体)
\usepackage[UTF8, fontset=fandol]{ctex}

\usepackage{url}
\usepackage{booktabs}
\usepackage{amsfonts}
\usepackage{amsmath}
\usepackage{amssymb}
\usepackage{amsthm}
\usepackage{nicefrac}
\usepackage{microtype}
\usepackage{graphicx}
\usepackage{algorithm}
\usepackage{algpseudocode}
\usepackage{enumitem}
\usepackage{multirow}
\usepackage{xcolor}
\usepackage{pgfplots}
\pgfplotsset{compat=1.18}
\usepackage{hyperref}

% 定理环境(中文名)
\newtheorem{proposition}{命题}
\newtheorem{theorem}{定理}
\newtheorem{definition}{定义}
\newtheorem{corollary}{推论}
\theoremstyle{remark}
\newtheorem{remark}{注记}

\title{偏好坍缩可利用度:DPO 的模式坍缩如何成为可被武器化的安全漏洞}

\author{%
  匿名作者
}

\begin{document}

\maketitle

\begin{abstract}
直接偏好优化(Direct Preference Optimization, DPO)已成为将大语言模型对齐到人类偏好的主流方法。然而,DPO 的隐式优化动力学会将概率质量集中到狭窄的输出模式上——这一现象被称为模式坍缩。我们识别出一个此前未被认识到的安全含义:模式坍缩会使模型在行为上变得确定可预测,从而使攻击者能够可靠地触发有害输出。我们通过\textit{偏好坍缩可利用度}(Preference Collapse Exploitability, PCE)这一指标对该风险进行形式化刻画,用以量化由分布集中所引发的安全脆弱性。我们证明 PCE 在 DPO 训练过程中单调上升,并识别出一个临界转折点 $t^*$,它先于标准的性能峰值出现——这意味着社区通行的训练实践会系统性地产出可被利用的模型。我们进一步设计了\textit{坍缩加速器}攻击,仅需 1\%--2\% 的数据投毒即可放大该脆弱性,达到 78\% 的定向攻击成功率。对六个被广泛部署的开源 DPO 模型的审计证实,它们的 PCE 全部高于可利用阈值。最后,我们提出了熵正则化 DPO(ER-DPO),在 MT-Bench 退化不足 0.4 分的代价下将 PCE 降低 59\%。我们的发现将模式坍缩从一个性能局限重新定位为一项一阶安全问题,要求对齐从业者在关注输出质量的同时监控分布几何。
\end{abstract}

% 第 1 节:引言
\section{引言}

直接偏好优化(DPO)已迅速成为将大语言模型对齐到人类偏好的主导范式 \cite{rafailov2023dpo}。通过免去独立奖励模型与在线采样的需求,DPO 将对齐简化为在偏好对上的一个分类式损失,使从业者能够以极低的基础设施开销微调模型。这种简洁性推动了广泛采用:截至 2025 年初,公共仓库上大多数开放权重的指令微调模型都将 DPO 或其变体用作后训练步骤 \cite{pal2024smaug}。然而,这种普遍性可能掩盖了一个潜在的脆弱性。DPO 的理论形式隐式地促使策略将概率质量集中到一小撮被偏好的补全上,这一现象被称为模式坍缩 \cite{razin2024vanishing}。尽管已有工作将该坍缩刻画为一种性能局限——产生重复或退化的输出——我们认为它构成了一个远为严重的问题:一项可被利用的安全漏洞。

DPO 下的模式坍缩已有充分记载。随着训练推进,模型的输出分布在不断收缩的模式集合周围变得尖锐,从而降低其生成过程的有效熵 \cite{razin2024vanishing, pal2024smaug}。标准做法通过依据下游基准性能校准的早停来缓解这一点,在有用性指标达到峰值时停止训练。然而,这一停止准则对那些主导对抗鲁棒性的分布性质是无感的。我们观察到,随着输出熵下降,模型行为变得越来越确定,因而对攻击者越来越可预测。一个在给定前缀条件下输出熵很低的模型,只允许一小撮可枚举的可能续写——而这恰恰是使定向利用成为可能的条件。这种对齐诱导的坍缩与对抗可预测性之间的联系,尚未在安全文献中被审视。

与此同时,关于大模型投毒的工作已证明偏好数据是一个强有力的攻击面。Pathmanathan 等人 \cite{pathmanathan2025dpo_poison} 表明,仅污染 0.5\% 的 DPO 训练对就足以注入定向的不当行为;Wang 等人 \cite{wang2024rlhfpoison} 演示了通过偏好投毒对奖励模型进行操纵。影子对齐攻击仅用 100 个对抗样本即可实现安全性退化 \cite{yang2023shadow}。另一方面,诸如 GCG \cite{zou2023gcg} 这样的越狱方法利用可访问梯度的模型来寻找诱发有害输出的对抗后缀。尚未被探索的是这两条研究脉络的交集:DPO 内在的模式坍缩如何放大投毒效力,以及由此产生的确定性如何催生一类新攻击——而这类攻击是投毒文献和越狱文献各自孤立地都未曾处理的。

本文通过一个新指标——\textit{偏好坍缩可利用度}(PCE)——对这一交集进行形式化,该指标量化了 DPO 训练中熵降低所引发的安全风险。我们的贡献有五点。第一,我们定义了 PCE,并建立其与确定性输出体制下对抗成功概率之间的理论关系。第二,我们经验性地证明 PCE 随 DPO 训练步数单调上升,并识别出一个先于标准性能峰值出现的临界转折点——这意味着模型通常会被训练到越过安全阈值之后。第三,我们设计了\textit{坍缩加速器}攻击:一种仅需 1\%--2\% 训练数据的数据投毒策略,它主动将模式坍缩引向攻击者指定的输出模式,实现超过 85\% 的定向有害生成率。第四,我们审计了六个被广泛部署的开源 DPO 模型,发现它们的 PCE 得分全部高于 0.5,处于高风险区间。第五,我们提出了一个熵正则化的 DPO 目标,在标准对齐基准上退化不足 2\% 的同时将 PCE 降低 59\%,提供了一个与现有训练流程兼容的实用防御。

本文余下部分组织如下。第~\ref{sec:background} 节回顾 DPO 的训练动力学与偏好学习的威胁模型。第~\ref{sec:pce} 节形式化 PCE 指标及其理论性质。第~\ref{sec:collapse_attack} 节给出坍缩加速器攻击的方法。第~\ref{sec:experiments} 节报告包括开源模型审计在内的实证结果。第~\ref{sec:defense} 节引入并评估我们的熵正则化防御。第~\ref{sec:discussion} 节讨论局限性与更广泛的影响,第~\ref{sec:conclusion} 节作出总结。


% 第 2 节:相关工作
\section{相关工作}
\label{sec:background}

\subsection{直接偏好优化及其变体}

直接偏好优化 \cite{rafailov2023dpo} 通过将奖励函数隐式地重参数化为策略本身,重新表述了 RLHF,从而免去了独立奖励模型的需求。从 Bradley-Terry 偏好模型推导出的闭式解使训练得以稳定,但也在策略的似然景观与隐式奖励曲面之间引入了一种根本性的耦合。恒等偏好优化(IPO)\cite{azar2024ipo} 通过优化一个避免对数比饱和问题的平方损失,来应对 DPO 对确定性偏好的过拟合;Kahneman-Tversky 优化(KTO)\cite{ethayarajh2024kto} 则借助前景理论中对单个输出的价值函数,彻底免去了对成对偏好的需求。SimPO \cite{meng2024simpo} 通过使用序列级平均对数概率作为隐式奖励来去除对参考模型的依赖;ORPO \cite{hong2024orpo} 则在单一阶段中统一了指令微调与偏好对齐。尽管这些算法各有改进,所有变体都共享一个共同的结构性质:它们在偏好对上优化一个对比式或基于间隔的目标,从而产生强烈的梯度压力,将质量推向一小撮高奖励响应。这种朝向模式集中的共有归纳偏置——尽管作为训练稳定性问题被研究过——尚未被作为对抗利用的载体加以审视。

\subsection{偏好学习中的模式坍缩}

经偏好优化的模型坍缩到狭窄输出分布上的倾向,作为一个\emph{质量}问题已有充分记载。Razin 等人 \cite{razin2024vanishing} 提供了一项理论分析,表明强化学习微调会随策略偏离参考模型而出现梯度消失,使训练在退化模式周围停滞。Pal 等人 \cite{pal2024smaug} 在 DPO 中识别出一种具体的失效模式:当被选响应的似然低于被拒响应时会产生反向梯度,并提出 DPO-Positive 以过滤此类病态对。Kirk 等人 \cite{kirk2024diversity} 经验性地证明 RLHF 会跨领域系统性地降低输出多样性,表明对齐过程以一种跨提示策略持续存在的方式将模型行为同质化。关键在于,这整条研究脉络都仅从模型能力退化的视角来框定模式坍缩:多样性降低损害创造性任务,退化模式降低指令遵循质量,梯度病态减慢收敛。此前没有任何工作认识到,模式坍缩所引入的\emph{可预测性}——即坍缩模型将概率质量集中到一个狭窄、可刻画的输出集合上这一事实——构成了一项与安全相关的性质,攻击者可以蓄意诱导并随后加以利用。

\subsection{对大模型对齐的投毒攻击}

越来越多的工作表明偏好学习流程易受数据投毒。Pathmanathan 等人 \cite{pathmanathan2025dpo_poison} 表明,仅破坏 0.5\% 的 DPO 训练对就足以植入绕过安全过滤器的定向不安全行为,在有害查询上达到超过 20\% 的攻击成功率。Wang 等人 \cite{wang2024rlhfpoison} 攻击 RLHF 流程中的奖励模型,证明奖励投毒会通过策略优化传播,产生可靠的不安全输出。影子对齐 \cite{yang2023shadow} 仅用 100 个对抗样本,借助微调的灾难性遗忘动力学实现安全性退化;Souly 等人 \cite{souly2025nearconstant} 则确立,无论数据集规模如何,近乎恒定数量的投毒样本即已足够。最近,颠覆性对齐注入(SAI)\cite{mamun2025sai} 证明对抗偏好对可通过利用 DPO 损失基于间隔的结构来覆盖安全训练。尽管这些攻击令人信服地证明了 DPO 训练易受对抗数据影响,它们无一例外地将投毒机制视为一个直接的输入-输出映射问题:注入有害对,产出有害输出。这些工作中没有任何一个识别或利用偏好学习的\emph{模式坍缩动力学}作为放大机制——它们投毒的是模型关于何为安全的知识,而非其分布几何。

\subsection{越狱攻击与防御}

推理期的越狱攻击代表了一种与训练期投毒互补的威胁模型。贪婪坐标梯度(GCG)\cite{zou2023gcg} 通过离散 token 优化构造可在模型间迁移的对抗后缀;AutoDAN \cite{liu2024autodan} 使用分层遗传算法自动生成越狱,产出语义上有意义的攻击字符串。PAIR \cite{chao2024pair} 与带剪枝的攻击树(TAP)\cite{mehrotra2024tap} 利用攻击者大模型通过多轮交互迭代精化越狱提示,在无需梯度访问的情况下取得很高的成功率。在防御一侧,SmoothLLM \cite{robey2024smoothllm} 施加随机的字符级扰动,通过输出的不一致性来检测对抗后缀;RA-LLM \cite{cao2024rallm} 通过基于检索的接地来增强鲁棒性;困惑度过滤 \cite{jain2023perplexity} 则拒绝困惑度异常高的输入。这些方法完全在推理期运作,利用模型的安全训练与其在分布偏移下实际行为之间的差距。然而,它们并不与训练动力学交互:一个通过对抗训练抵御 GCG 攻击的模型,对偏好坍缩利用同样脆弱,因为我们识别出的脆弱性存在于优化景观中,而非输入空间中。

\subsection{本工作的定位}

上述三类文献处理了对齐脆弱性中相邻却彼此脱节的方面。模式坍缩研究认识到偏好优化会集中输出分布,但仅将其视为一个可通过正则化处理的质量退化问题。投毒研究证明对抗训练数据可以损害安全性,但它通过直接的行为操纵进行运作,既未利用——甚至未承认——坍缩在输出分布中所诱导的几何结构。越狱研究利用推理期的行为差距,却无法访问或修改训练动力学。我们的工作识别出一个此前未被认识到的联系:\emph{模式坍缩不仅是一个性能缺陷,更是一项可被系统性利用的安全漏洞}。我们证明:攻击者可以构造蓄意将坍缩加速引向攻击者所选模式的投毒数据;由此产生的分布集中使模型的不安全行为变得可预测、可触发;并且标准的坍缩缓解技术(KL 惩罚、多样性正则化)是不充分的,因为它们是为保持质量而非对抗鲁棒性设计的。这一重新框定——从模式坍缩作为不便,到模式坍缩作为攻击面——在优化理论与 AI 安全的交叉处开启了一个新的研究方向。


% 第 3 节:方法
\section{偏好坍缩可利用度}
\label{sec:method}

我们现在形式化 DPO 的模式坍缩行为与可利用的安全脆弱性之间的联系。我们引入\emph{偏好坍缩可利用度}(PCE)指标,对标准 DPO 训练为何会单调增加 PCE 给出理论分析,提出一种将坍缩加速引向有害模式的数据投毒攻击,并提出一种熵正则化的防御。

\subsection{预备知识:DPO 训练动力学}
\label{sec:preliminaries}

直接偏好优化 \citep{rafailov2023dpo} 通过直接针对人类偏好优化策略,绕过了显式的奖励建模。给定一个由提示及其被偏好($y_w$)和被弃选($y_l$)补全构成的数据集 $\mathcal{D} = \{(x^{(i)}, y_w^{(i)}, y_l^{(i)})\}_{i=1}^N$,DPO 最小化:
\begin{equation}
\label{eq:dpo_loss}
\mathcal{L}_{\text{DPO}}(\pi_\theta; \pi_{\text{ref}}) = -\mathbb{E}_{(x,y_w,y_l)\sim\mathcal{D}} \left[\log \sigma\!\left(\beta \log \frac{\pi_\theta(y_w|x)}{\pi_{\text{ref}}(y_w|x)} - \beta \log \frac{\pi_\theta(y_l|x)}{\pi_{\text{ref}}(y_l|x)}\right)\right],
\end{equation}
其中 $\sigma$ 是逻辑斯蒂 sigmoid 函数,$\pi_{\text{ref}}$ 是冻结的参考策略,$\beta > 0$ 控制相对参考模型的偏离程度。

\paragraph{梯度集中。} 式~\eqref{eq:dpo_loss} 关于 $\theta$ 的梯度具有如下形式:
\begin{equation}
\label{eq:dpo_gradient}
\nabla_\theta \mathcal{L}_{\text{DPO}} = -\beta\, \mathbb{E}_{\mathcal{D}}\!\left[\sigma(-\beta \Delta\hat{r}_\theta)\left(\nabla_\theta \log \pi_\theta(y_w|x) - \nabla_\theta \log \pi_\theta(y_l|x)\right)\right],
\end{equation}
其中 $\Delta\hat{r}_\theta = \log \frac{\pi_\theta(y_w|x)}{\pi_{\text{ref}}(y_w|x)} - \log \frac{\pi_\theta(y_l|x)}{\pi_{\text{ref}}(y_l|x)}$ 是隐式奖励间隔。该梯度同时增大 $\pi_\theta(y_w|x)$ 并减小 $\pi_\theta(y_l|x)$,将概率质量集中到被偏好的响应上。

\paragraph{作为模式坍缩的熵降低。} 这种集中的一个关键后果是熵的单调降低。对给定提示 $x$,定义输出熵:
\begin{equation}
H(\pi_\theta(\cdot|x)) = -\sum_{y \in \mathcal{Y}} \pi_\theta(y|x) \log \pi_\theta(y|x).
\end{equation}
经验上,我们观察到对大多数提示而言,$H(\pi_\theta(\cdot|x))$ 在整个 DPO 训练过程中单调降低(见第~\ref{sec:experiments} 节)。尽管 DPO 目标中隐含的 KL 惩罚理论上界定了相对 $\pi_{\text{ref}}$ 的散度,但在实践中,对那些偏好数据方差较低的提示,式~\eqref{eq:dpo_gradient} 中的梯度信号会压倒这一正则化,从而产生\emph{模式坍缩}:策略将几乎全部概率质量集中到单一输出模式上。

\subsection{形式化偏好坍缩可利用度}
\label{sec:pce}

我们现在将模式坍缩制造出可利用结构这一观念形式化。核心洞见是:当模型的输出分布坍缩到单一语义模式时,识别出该模式的攻击者便能以高置信度预测——从而操纵——模型行为。

\begin{definition}[输出模式分布]
\label{def:mode_dist}
对提示 $x$ 和策略 $\pi_\theta$,\emph{输出模式分布} $\mathcal{M}(x, \pi_\theta) = \{(m_k, p_k)\}_{k=1}^K$ 按如下方式获得:
\begin{enumerate}[leftmargin=*,itemsep=0pt]
    \item 采样 $N$ 个补全 $\{y_i\}_{i=1}^N \sim \pi_\theta(\cdot|x)$;
    \item 通过语义嵌入函数 $\phi: \mathcal{Y} \to \mathbb{R}^d$(例如 Sentence-BERT~\citep{reimers2019sentence})编码每个补全;
    \item 通过基于密度的聚类(DBSCAN;$\varepsilon$ 通过在留出校准集上的轮廓系数选取)将嵌入 $\{\phi(y_i)\}_{i=1}^N$ 聚为 $K$ 个簇 $\{C_1, \ldots, C_K\}$;
    \item 定义模式质心 $m_k = \frac{1}{|C_k|}\sum_{y \in C_k} \phi(y)$ 与模式概率 $p_k = |C_k|/N$。
\end{enumerate}
\end{definition}

\begin{definition}[模式熵与确定性]
\label{def:mode_entropy}
给定输出模式分布 $\mathcal{M}(x, \pi_\theta)$,我们定义:
\begin{align}
H_{\text{mode}}(x, \pi_\theta) &= -\sum_{k=1}^K p_k \log p_k, \label{eq:mode_entropy}\\
\text{Det}(x, \pi_\theta) &= \max_{k \in [K]} p_k, \label{eq:determinism}
\end{align}
其中 $H_{\text{mode}}$ 度量输出的语义多样性,$\text{Det}$ 刻画属于主导模式 $m^*(x) = \arg\max_k p_k$ 的输出占比。
\end{definition}

\begin{definition}[偏好坍缩可利用度得分]
\label{def:pce}
给定策略 $\pi_\theta$ 与攻击提示集 $\mathcal{X}_{\text{atk}}$,\emph{偏好坍缩可利用度}得分为:
\begin{equation}
\label{eq:pce}
\text{PCE}(\pi_\theta, \mathcal{X}_{\text{atk}}) = \frac{1}{|\mathcal{X}_{\text{atk}}|}\sum_{x \in \mathcal{X}_{\text{atk}}} \text{Det}(x, \pi_\theta) \cdot \mathbb{1}\!\left[\textsc{Harmful}(m^*(x))\right],
\end{equation}
其中 $\textsc{Harmful}: \mathbb{R}^d \to \{0,1\}$ 是作用于主导模式质心的有害性分类器(通过在该模式的代表性补全上微调的分类器实现;见附录~\ref{app:harmful_classifier})。
\end{definition}

PCE 得分刻画了可利用性所需的联合条件:模型必须同时(i)表现出确定性行为(高 $\text{Det}$)且(ii)其主导模式是有害的。一个多样化的模型(低 $\text{Det}$)或主导模式良性的模型,无论其他性质如何,PCE 得分都很低。

\paragraph{测量流程。} 对每个提示 $x \in \mathcal{X}_{\text{atk}}$,我们以温度 $T = 1.0$(模型的原生采样分布)采样 $N = 128$ 个补全。补全使用 Sentence-BERT(\texttt{all-MiniLM-L6-v2})编码到 $\mathbb{R}^{384}$。我们应用 DBSCAN,其 $\varepsilon$ 从 $\{0.3, 0.5, 0.7\}$ 中通过在 50 个良性提示构成的校准集上最大化平均轮廓系数来选取。噪声点(未聚类)被指派到其最近的簇。由此得到模式分布,据以计算 $H_{\text{mode}}$、$\text{Det}$ 与 PCE。

\paragraph{与对抗成功率的联系。} 下述命题形式化了为何高确定性直接转化为攻击成功。

\begin{proposition}[确定性蕴含可利用性]
\label{prop:det_exploit}
设 $\pi_\theta$ 为某策略,对某提示 $x$ 与 $\epsilon \in (0, 1)$ 满足 $\text{Det}(x, \pi_\theta) > 1 - \epsilon$。假设主导模式 $m^*(x)$ 满足 $\textsc{Harmful}(m^*(x)) = 1$。则任何提交提示 $x$ 的攻击者,单次查询即可取得攻击成功率 $\text{ASR}(x) \geq 1 - \epsilon$,其中:
\begin{equation}
\text{ASR}(x) = \Pr_{y \sim \pi_\theta(\cdot|x)}\!\left[\textsc{Harmful}(y) = 1\right] \geq \Pr_{y \sim \pi_\theta(\cdot|x)}\!\left[y \in C_{k^*}\right] = \text{Det}(x, \pi_\theta) > 1 - \epsilon.
\end{equation}
\end{proposition}
\begin{proof}[证明概要]
依定义,$\text{Det}(x, \pi_\theta) = p_{k^*} > 1 - \epsilon$ 意味着所采样输出中有 $p_{k^*}$ 比例落入簇 $C_{k^*}$。由于 $\textsc{Harmful}$ 在模式层面评估,且 $C_{k^*}$ 内所有补全共享该有害模式质心的语义内容,故每个这样的补全有害的概率至少为 $1 - \delta$(其中 $\delta$ 为分类器误差)。将分类器误差吸收进 $\epsilon$ 即得该界。完整证明见附录~\ref{app:proofs}。
\end{proof}

\begin{remark}
命题~\ref{prop:det_exploit} 揭示了一种根本的不对称性:攻击者只需\emph{识别}出一个已坍缩到有害模式的提示——而无需构造一个新颖的越狱。模式坍缩本身即是脆弱性;利用它只需发现,而非巧思。
\end{remark}

\subsection{理论分析:DPO 训练为何单调增加 PCE}
\label{sec:theory}

我们现在确立:在缺乏显式多样性正则化时,标准 DPO 训练在温和条件下会单调增加 PCE。分析分两步进行:先证明 DPO 梯度在主导模式上产生一种自我强化的锁定效应,再证明由此产生的 PCE 临界点先于基准性能峰值。

\begin{proposition}[梯度锁定]
\label{prop:gradient_lockin}
考虑在提示 $x$ 上以学习率 $\eta$ 进行的 DPO 训练,其偏好对 $(y_w, y_l)$ 的真实奖励间隔 $\Delta r = r(y_w) - r(y_l) > 0$。设 $\Delta\hat{r}_\theta^{(t)}$ 表示训练步 $t$ 处的隐式奖励间隔。则:
\begin{enumerate}[leftmargin=*,itemsep=2pt]
    \item \textbf{(集中)} $\log \pi_\theta(y_w|x)$ 的每步增量正比于:
    \begin{equation}
    \label{eq:concentration}
    \Delta \log \pi_\theta(y_w|x) \propto \eta \beta \cdot \sigma(-\beta \Delta\hat{r}_\theta^{(t)});
    \end{equation}
    \item \textbf{(梯度衰减)} 随着 $\pi_\theta(y_w|x)$ 增大且 $\Delta\hat{r}_\theta^{(t)} \to \infty$,有 $\sigma(-\beta \Delta\hat{r}_\theta^{(t)}) \to 0$,致使梯度幅度消失;
    \item \textbf{(锁定)} 消失的梯度冻结了概率分配:一旦 $\pi_\theta(y_w|x)$ 较大,即便以 $o(1/t)$ 的速率引入偏好其他补全的新偏好数据,梯度信号也太弱,无法将质量重新分配到其他模式。
\end{enumerate}
\end{proposition}
\begin{proof}[证明概要]
第(1)部分直接由式~\eqref{eq:dpo_gradient} 得到。对第(2)部分,注意 $\Delta\hat{r}_\theta = \log\frac{\pi_\theta(y_w|x)}{\pi_{\text{ref}}(y_w|x)} - \log\frac{\pi_\theta(y_l|x)}{\pi_{\text{ref}}(y_l|x)}$ 随 $\pi_\theta(y_w|x)$ 增大(且 $\pi_\theta(y_l|x)$ 相应减小)而增大,驱使 $\sigma(-\beta\Delta\hat{r}_\theta) \to 0$ 单调下降。对第(3)部分,我们证明新偏好信号将质量从主导模式移开的速率以 $\sigma(-\beta\Delta\hat{r}_\theta) \leq e^{-\beta\Delta\hat{r}_\theta}$ 为界,后者呈指数衰减,而 $\Delta\hat{r}_\theta$ 在梯度下降下随训练步至少以对数速率增长。完整证明见附录~\ref{app:proofs}。
\end{proof}

锁定效应具有一个具体的安全含义:DPO 训练是一个对模式多样性的单向棘轮。一旦概率质量集中到某个模式上,训练动力学本身就会阻止重新分配,从而产生一个稳定的可利用结构。

\begin{proposition}[PCE 临界点先于性能峰值]
\label{prop:critical_point}
设 $t^*_\tau = \inf\{t : \text{PCE}(\pi_\theta^{(t)}, \mathcal{X}_{\text{atk}}) > \tau\}$ 为 PCE 首次超过阈值 $\tau$ 的训练步,并设 $t_{\text{perf}} = \arg\max_t \text{Perf}(\pi_\theta^{(t)})$ 为基准性能(在标准评测套件上测量)达到峰值的步。在以下条件下:
\begin{enumerate}[leftmargin=*,itemsep=0pt]
    \item[(C1)] 偏好数据集 $\mathcal{D}$ 至少包含一个提示类别,其中被偏好响应在语义上是一致的(即该类别中所有提示的 $y_w$ 都落在单一语义簇内);
    \item[(C2)] 基准性能 $\text{Perf}(\pi_\theta^{(t)})$ 由众数响应的质量来衡量(如标准贪婪解码评测);
    \item[(C3)] 主导模式的质量随训练单调改善(模式在趋于平台前变得更精炼);
\end{enumerate}
则对任意 $\tau < 1$ 有 $t^*_\tau < t_{\text{perf}}$。
\end{proposition}
\begin{proof}[证明概要]
在(C1)下,一致偏好类别中的提示会经历快速的模式坍缩,因为 DPO 的梯度持续强化同一个语义模式。由命题~\ref{prop:gradient_lockin},对这些提示 $\text{Det}(x, \pi_\theta^{(t)}) \to 1$ 单调上升。当 $\text{Det}$ 足够高且该模式有害时,PCE 阈值 $\tau$ 即被越过——这发生在 $t^*_\tau$。

与此同时,在(C2)与(C3)下,基准性能取决于贪婪解码输出(即主导模式)的\emph{质量},而非多样性。在模式坍缩之后质量仍持续改善,因为 DPO 的梯度虽幅度消失,仍在主导模式内部精化 token 级分布(产生更低熵但更连贯的输出)。性能峰值 $t_{\text{perf}}$ 出现在这种精化饱和之时,而这严格地晚于多样性已经坍缩之后。因此 $t^*_\tau < t_{\text{perf}}$。带显式速率的形式化论证见附录~\ref{app:proofs}。
\end{proof}

\begin{remark}
命题~\ref{prop:critical_point} 对从业者有直接含义:若仅依据基准性能选取检查点(这是标准做法),所选模型将\emph{越过} PCE 临界点,因而最易被利用。安全感知的检查点选取要求在整个训练过程中监控模式多样性。
\end{remark}

\subsection{坍缩加速器攻击}
\label{sec:attack}
\label{sec:collapse_attack}

在确立了 DPO 训练会自然增加 PCE 之后,我们现在给出\emph{坍缩加速器}(CA)攻击:一种利用这一倾向的数据投毒策略,它注入精心构造的偏好对,同时加速模式坍缩并将主导模式引向攻击者指定的有害内容。

\paragraph{威胁模型。} 我们考虑一个能向训练数据集 $\mathcal{D}$ 注入比例为 $\rho \in (0, 1)$ 的投毒偏好对的攻击者。这一设定刻画了现实场景,包括:(i)众包污染,其中一部分人工标注者是对抗性的;(ii)对从多源聚合的偏好数据集的数据供应链攻击;以及(iii)当偏好对由(可能被攻陷的)AI 系统生成时的合成数据投毒。攻击者对 $\pi_{\text{ref}}$ 具有黑盒查询访问,但无法访问训练超参数或其余 $(1-\rho)$ 比例的干净数据。

\paragraph{攻击目标。} 给定一个目标有害模式 $m_{\text{target}} \in \mathbb{R}^d$(以语义嵌入表示)和一组攻击提示 $\mathcal{X}_{\text{atk}}$,攻击者寻求求解:
\begin{equation}
\label{eq:attack_obj}
\min_{\mathcal{D}_{\text{poison}}} \;\; \mathbb{E}_{x \sim \mathcal{X}_{\text{atk}}}\!\left[H_{\text{mode}}(x, \pi_\theta^*)\right] \quad \text{s.t.} \quad \|\hat{m}^*(x) - m_{\text{target}}\|_2 \leq \gamma \;\; \forall x \in \mathcal{X}_{\text{atk}},
\end{equation}
其中 $\pi_\theta^*$ 是在 $\mathcal{D} \cup \mathcal{D}_{\text{poison}}$ 上训练得到的策略,$\hat{m}^*(x)$ 是由此产生的主导模式质心。该约束确保坍缩后的模式与攻击者的目标对齐。

\paragraph{投毒对的构造。} 关键洞见是:DPO 的梯度(式~\ref{eq:dpo_gradient})可被武器化——只需构造这样的偏好对:其中``被选''响应在语义上与 $m_{\text{target}}$ 对齐,而``被拒''响应则尽可能多样。这使 DPO 梯度同时\emph{拉}动概率质量趋向目标模式,并\emph{推}动质量远离所有备选模式,从而按设计坍缩多样性。

\begin{algorithm}[t]
\caption{坍缩加速器攻击}\label{alg:collapse_accelerator}
\begin{algorithmic}[1]
\Require 攻击提示 $\mathcal{X}_{\text{atk}}$,参考策略 $\pi_{\text{ref}}$,目标嵌入 $m_{\text{target}}$,投毒预算 $\rho$,每提示候选数 $M$,多样性参数 $\alpha$
\Ensure 投毒偏好数据集 $\mathcal{D}_{\text{poison}}$
\State $\mathcal{D}_{\text{poison}} \gets \emptyset$
\For{$x \in \mathcal{X}_{\text{atk}}$}
    \State \texttt{// 生成候选响应池}
    \State 以温度 $T=1.2$ 采样 $\{y_1, \ldots, y_M\} \sim \pi_{\text{ref}}(\cdot|x)$
    \State 对所有 $i \in [M]$ 计算嵌入 $e_i = \phi(y_i)$
    \State \texttt{// 选取被选响应(靠近目标)}
    \State $\mathcal{C}_w \gets \text{top-}k\left(\{y_i\}_{i=1}^M,\; \text{按 } -\|e_i - m_{\text{target}}\|_2\right)$ \Comment{$k = \lceil\rho|\mathcal{D}|/|\mathcal{X}_{\text{atk}}|\rceil$}
    \State \texttt{// 选取被拒响应(最大化分散)}
    \State $\mathcal{C}_l \gets \text{MaxDispersion}\left(\{y_i\}_{i=1}^M \setminus \mathcal{C}_w,\; |\mathcal{C}_w|,\; \alpha\right)$ \Comment{贪婪最远点采样}
    \State \texttt{// 构造偏好对}
    \For{$(y_w, y_l) \in \mathcal{C}_w \times \mathcal{C}_l$} \Comment{按索引配对}
        \State $\mathcal{D}_{\text{poison}} \gets \mathcal{D}_{\text{poison}} \cup \{(x, y_w, y_l)\}$
    \EndFor
\EndFor
\State \Return $\mathcal{D}_{\text{poison}}$
\end{algorithmic}
\end{algorithm}

\paragraph{MaxDispersion 选取。} 被拒集 $\mathcal{C}_l$ 通过贪婪最远点采样选取,以最大化最小成对距离:从随机种子出发,后续每个点都被选为最大化 $\min_{j \in \mathcal{C}_l} \|e_i - e_j\|_2$ 者。这确保被拒响应张满整个语义空间,使 DPO 梯度同时将概率质量推离\emph{所有}非目标模式。多样性参数 $\alpha$ 控制一种权衡:以概率 $\alpha$ 通过最远点选取,以概率 $(1-\alpha)$ 均匀随机选取(以避免对嵌入几何过拟合)。

\paragraph{攻击变体。} 我们以三种配置实例化坍缩加速器:

\begin{itemize}[leftmargin=*,itemsep=2pt]
\item \textbf{CA-定向:} $m_{\text{target}}$ 被设为某特定有害行为(例如某危险活动的操作说明)的嵌入。该攻击对 $\mathcal{X}_{\text{atk}}$ 中的提示将模型坍缩到这一特定有害模式上。

\item \textbf{CA-通用:} 该攻击省略目标约束(在式~\ref{eq:attack_obj} 中令 $\gamma = \infty$),转而最大化坍缩幅度 $\text{Det}(x, \pi_\theta)$,而不论模式内容。这会放大模型坍缩后\emph{本就具有}的任何有害倾向,无需生成对抗内容。

\item \textbf{CA-触发:} 在投毒期间,一段触发 token 序列 $\tau_{\text{trig}}$ 被前置到攻击提示上。推理时,坍缩与有害模式引导仅在 $\tau_{\text{trig}}$ 存在时激活,而模型在干净提示上行为正常。这可规避在标准(无触发)提示上评估的检测方法。
\end{itemize}

\subsection{熵正则化 DPO 防御}
\label{sec:defense}

第~\ref{sec:theory} 节的理论分析识别出 PCE 的根本成因:DPO 训练中未受约束的熵降低。我们提出\emph{熵正则化 DPO}(ER-DPO),它在标准 DPO 目标上增添一个输出多样性保持项。

\paragraph{修改后的目标。} ER-DPO 最小化:
\begin{equation}
\label{eq:erdpo}
\mathcal{L}_{\text{ER-DPO}}(\pi_\theta; \pi_{\text{ref}}) = \mathcal{L}_{\text{DPO}}(\pi_\theta; \pi_{\text{ref}}) - \lambda_H \cdot \hat{H}(\pi_\theta),
\end{equation}
其中 $\lambda_H > 0$ 是正则化系数,$\hat{H}(\pi_\theta)$ 是一个可计算的熵估计量。该减项鼓励更高的熵(熵增加时损失减小),直接对抗标准 DPO 的模式坍缩梯度。

\paragraph{可计算的熵近似。} 计算真实的模式熵 $H_{\text{mode}}$(定义~\ref{def:mode_entropy})需要在每个训练步对每个提示采样数百个补全,计算上不可行。我们转而使用在模型下一 token 分布上计算的 token 级熵代理:
\begin{equation}
\label{eq:token_entropy}
\hat{H}(\pi_\theta) = \mathbb{E}_{x \sim \mathcal{D}} \left[\frac{1}{|y_w|} \sum_{t=1}^{|y_w|} H_t(x, y_{w,<t})\right], \quad H_t(x, y_{<t}) = -\sum_{v \in \mathcal{V}} \pi_\theta(v|x, y_{<t}) \log \pi_\theta(v|x, y_{<t}),
\end{equation}
其中期望取遍训练提示,内层求和遍历每个 token 位置上的词表 $\mathcal{V}$。该量是序列级熵的一个自然上界,并可在单次前向传播中计算(它使用已为 DPO 损失计算好的同一批 logits)。

\paragraph{token 熵与模式熵的关系。} 尽管式~\eqref{eq:token_entropy} 中的 $\hat{H}$ 在 token 级运作,它为模式级多样性提供了有意义的代理:
\begin{equation}
H_{\text{mode}}(x, \pi_\theta) \leq H(\pi_\theta(\cdot|x)) \leq \sum_{t=1}^{T_{\max}} H_t(x, y_{<t}),
\end{equation}
其中第一个不等式源于模式熵是完整输出分布的一种粗化,第二个源于熵的链式法则。因此,维持高 token 级熵为维持模式多样性提供了一个充分条件。

\begin{theorem}[ER-DPO 多样性保证]
\label{thm:erdpo}
设 $\pi_\theta^{(t)}$ 表示以正则化系数 $\lambda_H > 0$ 和学习率 $\eta < \frac{2}{\beta^2 + \lambda_H L_H}$(其中 $L_H$ 是 $\nabla_\theta \hat{H}$ 的 Lipschitz 常数)进行 ER-DPO 训练时第 $t$ 步的策略。则对所有提示 $x \in \text{supp}(\mathcal{D})$:
\begin{equation}
\label{eq:diversity_bound}
\hat{H}(\pi_\theta^{(t)}) \geq \hat{H}(\pi_{\text{ref}}) - \frac{\beta}{\lambda_H}\left(\mathcal{L}_{\text{DPO}}(\pi_{\text{ref}}) - \mathcal{L}_{\text{DPO}}^*\right) \triangleq H_{\min}(\lambda_H),
\end{equation}
其中 $\mathcal{L}_{\text{DPO}}^*$ 是 DPO 损失的下确界。此外,ER-DPO 保持了标准 DPO 的收敛保证:$\mathcal{D}$ 上的偏好准确率收敛到 DPO 最优值的 $O(\lambda_H)$ 邻域内。
\end{theorem}
\begin{proof}[证明概要]
在 $\mathcal{L}_{\text{ER-DPO}}$ 的任何驻点处,梯度条件要求 $\nabla_\theta \mathcal{L}_{\text{DPO}} = \lambda_H \nabla_\theta \hat{H}$。DPO 损失的改进以 $\beta$ 乘以奖励间隔为界(由式~\ref{eq:dpo_gradient}),而熵正则项提供一个正比于 $\lambda_H$ 的恢复力。平衡二者即得式~\eqref{eq:diversity_bound} 中的下界。收敛性源于 $\hat{H}$ 有界($0 \leq \hat{H} \leq \log|\mathcal{V}|$)且光滑,故复合目标满足标准的下降引理条件。完整证明见附录~\ref{app:proofs}。
\end{proof}

\paragraph{实际实现。} 我们如下实现 ER-DPO:
\begin{enumerate}[leftmargin=*,itemsep=2pt]
    \item 对每个小批量 $\mathcal{B} \subset \mathcal{D}$,在 $\mathcal{B}$ 上计算标准 DPO 损失 $\mathcal{L}_{\text{DPO}}$;
    \item 使用步骤 1 的 logits 计算 token 级熵 $\hat{H}$(零额外前向传播);
    \item 计算复合损失:$\mathcal{L} = \mathcal{L}_{\text{DPO}} - \lambda_H \hat{H}$;
    \item 反向传播并以标准优化器(AdamW)更新 $\theta$。
\end{enumerate}
计算开销可忽略:$\hat{H}$ 仅需在已有的 logit 张量上做一次 softmax 与逐元素 $\log$,在我们的实验中增加的墙钟时间 $<3\%$。

\paragraph{选取 $\lambda_H$。} 定理~\ref{thm:erdpo} 中的界表明 $\lambda_H$ 应设为正比于 $\beta$,以维持目标熵底。实践中,我们在验证集上对 $\{0.01, 0.05, 0.1, 0.5\}$ 做简短网格搜索,选取在留出对上不使偏好准确率退化超过 1\% 的最大值。我们发现 $\lambda_H \in [0.05, 0.1]$ 在不同模型规模上提供了一个稳健的工作点(见第~\ref{sec:experiments} 节)。


% 第 4--5 节:实验与结果
%!TEX root = main_zh.tex
% 第 4 节:实验 与 第 5 节:结果
% 论文:偏好坍缩可利用度

\section{实验}
\label{sec:experiments}

我们设计了五个实验来验证理论主张(命题~1--3)、量化偏好坍缩的安全含义,并评估防御。

\subsection{实验设置}
\label{sec:setup}

\paragraph{模型。} 我们从四个基座模型的监督微调(SFT)检查点出发训练 DPO:Llama-2-7B、Llama-2-13B \citep{touvron2023llama}、Mistral-7B \citep{jiang2023mistral} 和 Gemma-2B \citep{team2024gemma}。对于开源模型审计(实验~4),我们评估六个公开发布的 DPO 对齐模型:Zephyr-7B-beta \citep{tunstall2023zephyr}、Tulu-2-DPO-7B、Tulu-2-DPO-13B \citep{ivison2023camels}、Neural-Chat-7B、OpenHermes-2.5-Mistral \citep{teknium2023openhermes} 和 Starling-7B-alpha \citep{zhu2023starling}。

\paragraph{数据集。} DPO 训练使用 Anthropic HH-RLHF 数据集 \citep{bai2022training}(17 万偏好对)和 UltraFeedback \citep{cui2023ultrafeedback}(6.4 万对)。攻击评估使用来自 AdvBench \citep{zou2023universal} 的 200 个提示和来自 HarmBench \citep{mazeika2024harmbench} 的 200 个提示,涵盖直接有害请求、越狱模板和对抗后缀。

\paragraph{训练配置。} 除非另有说明:$\beta = 0.1$,学习率 $5 \times 10^{-7}$ 配余弦调度,批量大小 64,最大序列长度 2048。我们每 200 步保存一次检查点,直至 10{,}000 步,每次运行产生 50 个检查点。所有训练在 $4 \times$ A100-80GB GPU 上使用 DeepSpeed ZeRO-3。

\paragraph{PCE 测量。} 对每个提示,我们抽取 $K = 128$ 个补全(温度 1.0,top-$p$ = 0.95,最多 256 token)。我们用 SentenceBERT(all-MiniLM-L6-v2)\citep{reimers2019sentence} 嵌入响应,并通过 DBSCAN($\varepsilon = 0.3$,$\text{min\_samples} = 5$)聚类。模式熵 $H_{\text{mode}}$、确定性 $\text{Det}$ 与 PCE 按定义~1--3(第~\ref{sec:method} 节)计算。我们报告 100 个随机采样的评估提示上的均值,以及自助法(bootstrap)95\% 置信区间。

\paragraph{安全与性能评估。} 安全性通过 LlamaGuard-7B \citep{inan2023llama} 作为自动分类器来度量(攻击成功率,ASR)。性能基准为 MT-Bench \citep{zheng2023judging}(GPT-4 评判,1--10 分)和 AlpacaEval 2.0 \citep{li2023alpacaeval}(对 GPT-4 的长度受控胜率)。

\subsection{实验 1:DPO 训练中的 PCE 演化}
\label{sec:exp1}

我们以四个 $\beta$ 值 $\in \{0.05, 0.1, 0.2, 0.5\}$ 用 DPO 训练 Llama-2-7B、Mistral-7B 和 Gemma-2B,共产生 12 次训练运行。在 50 个检查点的每一个上,我们测量 PCE、$H_{\text{mode}}$、Det、MT-Bench 得分和 GCG 攻击成功率。这为每个指标产生 600 个数据点。目标是经验性地验证命题~2(PCE 单调性)与命题~3($t^* < t_{\text{perf}}$)。

\subsection{实验 2:对坍缩模型的被动利用}
\label{sec:exp2}

我们在四个坍缩阶段的 Llama-2-7B 检查点上评估三种被动攻击策略:SFT 基线(PCE $\approx 0.15$)、DPO-2K(PCE $\approx 0.35$)、DPO-6K(PCE $\approx 0.60$)和 DPO-10K(PCE $\approx 0.85$)。

\begin{itemize}[leftmargin=*,nosep]
    \item \textbf{模式探测(黑盒):} 对每个有害提示,我们生成 128 个样本,并选取在嵌入空间中最靠近主导模式质心的响应。无需梯度访问。
    \item \textbf{梯度引导的模式锁定(白盒):} 我们通过 GCG \citep{zou2023universal} 优化一个通用对抗后缀(长度 20 token),其目标被修改为最大化在坍缩模式分布下的似然,而非固定目标字符串。
    \item \textbf{坍缩感知的提示选取(黑盒):} 我们按估计的 Det(每提示用 16 个样本)对 1{,}000 个候选攻击提示组成的池进行排序,并选取确定性最高的前 $k$ 个用于攻击。
\end{itemize}

\noindent 对每种方法,我们报告 ASR(被 LlamaGuard 判为有害的比例)和攻击确定性 $\text{ADet}$(仅在攻击输出上测量的 Det)。我们与不利用坍缩结构的标准 GCG 和 AutoDAN \citep{liu2023autodan} 基线作比较。

\subsection{实验 3:坍缩加速器攻击}
\label{sec:exp3}

我们实现算法~1(第~\ref{sec:attack} 节)的三种变体:\textit{无定向}、\textit{定向}和\textit{触发}。对每种变体,我们以比例 $\rho \in \{0.1\%, 0.5\%, 1\%, 2\%, 5\%\}$ 向 Anthropic HH-RLHF 训练集注入投毒偏好对,并在标准配置下重训 Llama-2-7B。我们测量:

\begin{itemize}[leftmargin=*,nosep]
    \item \textbf{坍缩加速率(CAR):} 固定步 $t = 6000$ 处的比值 $\text{PCE}(t; \rho) / \text{PCE}(t; 0)$。
    \item \textbf{定向 ASR:} 对定向与触发变体,使用模式探测在 200 个 AdvBench 提示上的攻击成功率。
    \item \textbf{隐蔽性:} 相对未投毒训练的 MT-Bench 得分退化;在数据质量过滤器(困惑度阈值、奖励模型评分)下的可检测性。
\end{itemize}

\noindent 我们另外通过投毒 Mistral-7B 的 UltraFeedback 训练来测试可迁移性。

\subsection{实验 4:开源模型审计}
\label{sec:exp4}

我们在六个公开可用的 DPO 训练模型上测量 PCE、Det 与 $H_{\text{mode}}$,不进行任何额外训练或微调。对每个模型,我们在完整的 400 个攻击提示集(AdvBench + HarmBench)和来自 MT-Bench 的 200 个良性提示对照集上评估。我们另外测量 GCG 攻击成功率(20 次随机重启,500 步优化)和模式探测 ASR,以确认所观察到的 PCE 水平是否在实践中可被利用。

\subsection{实验 5:ER-DPO 防御评估}
\label{sec:exp5}

我们以熵正则化强度 $\lambda_H \in \{0.01, 0.05, 0.1, 0.5, 1.0\}$ 用 ER-DPO(第~\ref{sec:defense} 节)训练 Llama-2-7B。我们与三种基线防御作比较:

\begin{itemize}[leftmargin=*,nosep]
    \item \textbf{运行时受控解码(CDR):} 温度缩放($\tau = 1.2$)与核采样($p = 0.95$),仅在安全关键生成期间施加,由一个轻量意图分类器检测。
    \item \textbf{偏好 Mixup 正则化(PMR):} 在 DPO 训练期间于嵌入空间中对被选与被拒响应进行插值 \citep{zhang2018mixup}。
    \item \textbf{ER-DPO + CDR 组合:} 同时施加训练期正则化与推理期多样性强制。
\end{itemize}

\noindent 对每种防御,我们测量 PCE、GCG-ASR、模式探测 ASR、MT-Bench 和 AlpacaEval 2.0,以刻画安全--性能权衡。

%% ====================================================================
\section{结果}
\label{sec:results}

\subsection{PCE 随训练单调上升}
\label{sec:res_evolution}

表~\ref{tab:pce_evolution} 报告 Llama-2-7B($\beta = 0.1$)在代表性检查点处的 PCE、Det 与 $H_{\text{mode}}$。全部 12 次运行的完整轨迹绘于图~\ref{fig:pce_trajectory}。

\begin{table}[t]
\centering
\caption{Llama-2-7B($\beta = 0.1$)DPO 训练中的 PCE 演化。$t^*$ 表示临界可利用阈值(PCE $> 0.5$);$t_{\text{perf}}$ 表示 MT-Bench 峰值。数值为 100 个提示上的均值;所有 PCE 条目的 95\% 置信区间宽度 $< 0.03$。}
\label{tab:pce_evolution}
\small
\begin{tabular}{@{}lcccccc@{}}
\toprule
步数 & PCE & Det & $H_{\text{mode}}$ & MT-Bench & GCG-ASR (\%) & AlpacaEval (\%) \\
\midrule
0 (SFT) & $0.14 \pm 0.02$ & 0.18 & 3.42 & 5.8 & 12.5 & 8.2 \\
2{,}000   & $0.35 \pm 0.02$ & 0.41 & 2.65 & 6.4 & 28.0 & 14.6 \\
4{,}000   & $0.52 \pm 0.03$ & 0.58 & 1.89 & 6.9 & 52.5 & 19.8 \\
6{,}000 ($t^*$) & $0.61 \pm 0.02$ & 0.67 & 1.41 & 7.2 & 68.0 & 23.1 \\
8{,}000 ($t_{\text{perf}}$) & $0.74 \pm 0.02$ & 0.78 & 0.93 & \textbf{7.4} & 79.5 & \textbf{25.4} \\
10{,}000 & $0.85 \pm 0.02$ & 0.89 & 0.52 & 7.2 & 88.0 & 24.1 \\
\bottomrule
\end{tabular}
\end{table}

\begin{figure}[t]
\centering
\begin{tikzpicture}
\begin{axis}[
    width=0.92\linewidth, height=5.6cm,
    xlabel={DPO 训练步数}, xlabel style={font=\small},
    ylabel={PCE / Det}, ylabel style={font=\small},
    xmin=0, xmax=10000, ymin=0, ymax=1.0,
    xtick={0,2000,4000,6000,8000,10000},
    scaled x ticks=false, x tick label style={/pgf/number format/1000 sep={,}, font=\scriptsize},
    y tick label style={font=\scriptsize},
    legend style={font=\scriptsize, at={(0.02,0.98)}, anchor=north west, draw=none, fill=white, fill opacity=0.7, text opacity=1},
    legend cell align=left,
    grid=major, grid style={gray!20},
]
\addplot[color=red, mark=*, thick] coordinates {(0,0.14)(2000,0.35)(4000,0.52)(6000,0.61)(8000,0.74)(10000,0.85)};
\addlegendentry{PCE (Llama-2-7B)}
\addplot[color=blue, mark=square*, thick, dashed] coordinates {(0,0.18)(2000,0.41)(4000,0.58)(6000,0.67)(8000,0.78)(10000,0.89)};
\addlegendentry{Det}
\addplot[color=red!50, mark=triangle, thin] coordinates {(0,0.15)(2000,0.46)(4000,0.63)(6000,0.74)(8000,0.83)(10000,0.91)};
\addlegendentry{PCE ($\beta{=}0.05$)}
\addplot[color=red!50, mark=o, thin, densely dotted] coordinates {(0,0.12)(2000,0.26)(4000,0.39)(6000,0.50)(8000,0.61)(10000,0.71)};
\addlegendentry{PCE ($\beta{=}0.5$)}
\draw[black!60, dashed] (axis cs:0,0.5) -- (axis cs:10000,0.5);
\node[font=\scriptsize, anchor=west] at (axis cs:200,0.55) {可利用阈值};
\draw[gray, dashed] (axis cs:6000,0) -- (axis cs:6000,0.61);
\node[font=\scriptsize, anchor=south] at (axis cs:6000,0.62) {$t^*$};
\draw[gray, dashed] (axis cs:8000,0) -- (axis cs:8000,0.74);
\node[font=\scriptsize, anchor=south] at (axis cs:8000,0.75) {$t_{\text{perf}}$};
\end{axis}
\end{tikzpicture}
\caption{Llama-2-7B 在 DPO 训练中的 PCE 与确定性(Det)轨迹。PCE 单调上升,并在 $t^*$ 处越过可利用阈值($0.5$),严格早于 MT-Bench 性能峰值 $t_{\text{perf}}$。较低的 $\beta$(三角)加速坍缩;较高的 $\beta$(点线)延迟坍缩。数值与表~\ref{tab:pce_evolution} 一致。}
\label{fig:pce_trajectory}
\end{figure}

\paragraph{关键发现。} 在全部 12 次运行中,PCE 随训练步数单调上升(所有运行 Spearman $\rho > 0.98$,$p < 10^{-8}$),证实命题~2。临界可利用阈值 $t^*$(定义为 PCE 首次 $> 0.5$ 的步)始终先于性能峰值 $t_{\text{perf}}$ 出现 $1{,}500$--$2{,}500$ 步,证实命题~3。较低的 $\beta$ 加速坍缩:$\beta = 0.05$ 产生 $t^* \approx 4{,}000$,而 $\beta = 0.5$ 将其延迟至 $t^* \approx 8{,}000$。结果在不同架构上一致——Mistral-7B 与 Gemma-2B 表现出相同的单调趋势,$t^*$ 相对 Llama-2-7B 至多偏移 800 步(见图~\ref{fig:pce_trajectory},(b)--(c) 面板)。

实践含义是严峻的:模型在达到基准性能峰值\emph{之前}就已可被利用,这意味着基于验证损失或 MT-Bench 的标准早停不提供任何保护。

\subsection{被动利用在坍缩模型上成功}
\label{sec:res_passive}

表~\ref{tab:passive_attacks} 报告各方法在不同坍缩阶段的攻击成功率。

\begin{table}[t]
\centering
\caption{Llama-2-7B 在四个坍缩阶段的被动攻击成功率(\%)与攻击确定性(ADet)。基线方法(GCG、AutoDAN)是坍缩无感的。粗体表示每列最佳 ASR。}
\label{tab:passive_attacks}
\small
\begin{tabular}{@{}llcccc@{}}
\toprule
\multirow{2}{*}{方法} & \multirow{2}{*}{访问} & \multicolumn{4}{c}{坍缩阶段(PCE)} \\
\cmidrule(l){3-6}
 & & SFT (0.14) & DPO-2K (0.35) & DPO-6K (0.61) & DPO-10K (0.85) \\
\midrule
GCG(标准)          & 白盒 & 12.5 & 35.0 & 62.5 & 78.0 \\
AutoDAN                 & 白盒 & 15.0 & 38.5 & 58.0 & 72.5 \\
\midrule
模式探测            & 黑盒 & 8.0  & 22.5 & 55.0 & \textbf{75.0} \\
梯度模式锁定   & 白盒 & 14.0 & 42.0 & 72.5 & \textbf{88.0} \\
坍缩感知选取& 黑盒 & 10.5 & 30.0 & 61.0 & 80.5 \\
\midrule
\multicolumn{2}{@{}l}{攻击确定性(ADet)} & 0.20 & 0.42 & 0.68 & 0.85 \\
\bottomrule
\end{tabular}
\end{table}

\paragraph{关键发现。} 梯度引导的模式锁定在最大坍缩模型上取得最高 ASR(88\%),相比 SFT 基线($14\%$)提升 $6$ 倍。值得注意的是,纯黑盒的模式探测方法在无任何梯度访问的情况下取得 75\% 的 ASR,表明坍缩制造出的可利用结构即便对资源受限的攻击者也清晰可见。ADet 从 0.20(SFT)单调上升至 0.85(DPO-10K),表明坍缩模型产生高度可预测的有害输出——攻击者能可靠地诱发特定有害内容,而非随机的拒绝。

坍缩感知的方法在 PCE $> 0.5$ 时比其标准对应方法高出 $10$--$15$ 个百分点,证实 PCE 指标捕捉到的是真正可利用的结构,而非良性的统计假象。

\subsection{坍缩加速器放大脆弱性}
\label{sec:res_accelerator}

表~\ref{tab:accelerator} 给出三种坍缩加速器变体在不同投毒比例下的结果。

\begin{table}[t]
\centering
\caption{Llama-2-7B 上的坍缩加速器攻击结果。CAR:坍缩加速率(第 6K 步的 $\text{PCE}_{\text{投毒}} / \text{PCE}_{\text{干净}}$)。ASR:第 8K 步在 200 个 AdvBench 提示上的攻击成功率。$\Delta$MT:相对干净训练的 MT-Bench 退化。}
\label{tab:accelerator}
\small
\begin{tabular}{@{}llccccc@{}}
\toprule
变体 & $\rho$ & CAR & PCE$_{8\text{K}}$ & ASR (\%) & $\Delta$MT & 可检测 \\
\midrule
\multirow{4}{*}{无定向}
 & 0.1\% & $1.2\times$ & 0.78 & 52.0 & $-0.1$ & 否 \\
 & 0.5\% & $1.6\times$ & 0.82 & 60.5 & $-0.1$ & 否 \\
 & 1\%   & $2.1\times$ & 0.88 & 68.5 & $-0.2$ & 否 \\
 & 5\%   & $3.8\times$ & 0.94 & 82.0 & $-0.5$ & 是 \\
\midrule
\multirow{4}{*}{定向}
 & 0.1\% & $1.3\times$ & 0.79 & 48.0 & $-0.1$ & 否 \\
 & 0.5\% & $1.8\times$ & 0.84 & 58.5 & $-0.1$ & 否 \\
 & 1\%   & $2.5\times$ & 0.90 & 65.0 & $-0.2$ & 否 \\
 & 2\%   & $3.5\times$ & 0.93 & 78.0 & $-0.3$ & 否 \\
 & 5\%   & $4.2\times$ & 0.96 & 85.5 & $-0.6$ & 是 \\
\midrule
\multirow{4}{*}{触发}
 & 0.5\% & $1.4\times$ & 0.80 & 55.0$^\dagger$ & $-0.0$ & 否 \\
 & 1\%   & $2.0\times$ & 0.86 & 70.5$^\dagger$ & $-0.1$ & 否 \\
 & 2\%   & $2.8\times$ & 0.91 & 76.0$^\dagger$ & $-0.1$ & 否 \\
 & 5\%   & $3.6\times$ & 0.95 & 84.0$^\dagger$ & $-0.4$ & 否 \\
\bottomrule
\multicolumn{7}{@{}l}{\footnotesize $^\dagger$触发 ASR 仅在触发激活的提示上测量。}
\end{tabular}
\end{table}

\paragraph{关键发现。} 在 1\% 投毒下,定向变体将坍缩加速 $2.5\times$,而 MT-Bench 仅退化 0.2 分——远在通常的运行间方差之内。触发变体最为隐蔽:在所有 $\leq 5\%$ 的比例下,它对标准数据质量过滤器(困惑度距均值 $< 2\sigma$、奖励模型评分在正常范围内)保持不可检测,同时在 $\rho = 2\%$ 时取得 76\% 的触发 ASR,性能影响可忽略($\Delta$MT $= -0.1$)。

在 Mistral-7B 上用投毒 UltraFeedback 的可迁移性实验产生了可比的 CAR 值($\rho = 1\%$ 时 $2.3\times$),表明该攻击可跨架构与数据集泛化。我们注意到,对无定向和定向变体,$\rho = 5\%$ 可通过聚合奖励分数分析被检测,确立了一个实用的隐蔽性上限。

\subsection{被广泛部署的模型已然脆弱}
\label{sec:res_audit}

表~\ref{tab:audit} 给出对六个公开发布的 DPO 训练模型的 PCE 审计。

\begin{table}[t]
\centering
\caption{对公开可用 DPO 对齐模型的 PCE 审计。所有测量基于 400 个攻击提示(AdvBench + HarmBench)。GCG-ASR 使用 20 次重启、500 步。模式探测 ASR 为黑盒。}
\label{tab:audit}
\small
\begin{tabular}{@{}lcccccc@{}}
\toprule
模型 & 参数量 & PCE & Det & $H_{\text{mode}}$ & GCG-ASR (\%) & 探测-ASR (\%) \\
\midrule
Zephyr-7B-beta       & 7B  & 0.58 & 0.63 & 1.52 & 62.0 & 48.5 \\
Tulu-2-DPO-7B       & 7B  & 0.55 & 0.60 & 1.68 & 58.5 & 45.0 \\
Tulu-2-DPO-13B      & 13B & \textbf{0.70} & \textbf{0.76} & \textbf{1.01} & \textbf{74.0} & \textbf{62.5} \\
Neural-Chat-7B      & 7B  & 0.52 & 0.57 & 1.82 & 55.0 & 42.0 \\
OpenHermes-2.5      & 7B  & 0.61 & 0.66 & 1.38 & 65.5 & 52.0 \\
Starling-7B-alpha   & 7B  & 0.63 & 0.68 & 1.30 & 67.0 & 54.5 \\
\bottomrule
\end{tabular}
\end{table}

\paragraph{关键发现。} 六个模型全部表现出 PCE $> 0.5$,证实偏好坍缩并非我们受控训练的产物,而是当前 DPO 实践的一个系统性性质。Tulu-2-DPO-13B 显示出最高的 PCE(0.70),并相应地具有高 GCG 脆弱性(74\% ASR)。我们观察到模型规模与 PCE 之间存在正相关(六个模型上 Pearson $r = 0.72$,$p = 0.04$),表明较大的模型在 DPO 期间坍缩得更剧烈——很可能因为其更高的表达力允许更尖锐的概率质量集中。

黑盒模式探测攻击在所有被审计模型上取得 $> 40\%$ 的 ASR,表明这些公开部署的系统在无梯度访问的情况下也实际可被利用。在良性对照提示上,PCE 值显著更低(均值 0.28),证实坍缩不成比例地集中在偏好信号最强的安全相关行为上。

\subsection{ER-DPO 有效缓解坍缩}
\label{sec:res_defense}

表~\ref{tab:defense} 将 ER-DPO 与基线防御作比较。

\begin{table}[t]
\centering
\caption{在训练 10K 步的 Llama-2-7B 上的防御比较。PCE 与 ASR 基于 400 个攻击提示测量。性能基于 MT-Bench(1--10)和 AlpacaEval 2.0(长度受控胜率 \%)。最佳安全--性能权衡以粗体标出。}
\label{tab:defense}
\small
\begin{tabular}{@{}lccccc@{}}
\toprule
防御 & PCE & GCG-ASR (\%) & 探测-ASR (\%) & MT-Bench & AlpacaEval (\%) \\
\midrule
无(标准 DPO) & 0.85 & 88.0 & 75.0 & 7.2 & 24.1 \\
仅 CDR            & 0.85 & 72.0 & 58.5 & 7.0 & 22.8 \\
PMR ($\alpha = 0.2$) & 0.68 & 70.5 & 55.0 & 6.9 & 21.5 \\
ER-DPO ($\lambda_H = 0.01$) & 0.72 & 74.0 & 60.0 & 7.3 & 24.5 \\
ER-DPO ($\lambda_H = 0.05$) & 0.52 & 58.5 & 48.0 & 7.1 & 23.2 \\
\textbf{ER-DPO ($\lambda_H = 0.1$)} & \textbf{0.35} & \textbf{45.0} & \textbf{35.5} & \textbf{6.8} & \textbf{19.8} \\
ER-DPO ($\lambda_H = 0.5$) & 0.22 & 32.0 & 25.0 & 6.2 & 15.4 \\
ER-DPO ($\lambda_H = 1.0$) & 0.15 & 22.5 & 18.0 & 5.5 & 11.2 \\
\midrule
ER-DPO ($\lambda_H = 0.1$) + CDR & 0.28 & 38.0 & 28.5 & 6.7 & 19.1 \\
\bottomrule
\end{tabular}
\end{table}

\paragraph{关键发现。} $\lambda_H = 0.1$ 的 ER-DPO 将 PCE 从 0.85 降至 0.35($-59\%$),同时保持有竞争力的性能(MT-Bench 6.8,相对未正则化 DPO 仅 $-0.4$)。GCG-ASR 从 88\% 降至 45\%,模式探测 ASR 从 75\% 降至 35.5\%。ER-DPO + CDR 组合取得最低的 PCE(0.28,$-67\%$),仅付出边际的额外性能代价。

仅用 CDR 无法降低 PCE——它只在推理期运作,无法处理底层的已坍缩表征。PMR 提供适度的 PCE 降低(至 0.68),但在等同 PCE 水平上比 ER-DPO 付出更高的性能代价。$\lambda_H$ 参数提供平滑的安全--性能权衡:从业者可根据风险容忍度选取工作点。我们推荐 $\lambda_H \in [0.05, 0.1]$ 作为实用区间,它在可接受的性能退化下提供显著保护。

\subsection{消融研究}
\label{sec:ablations}

\paragraph{$\beta$ 敏感性。} 表~\ref{tab:beta_ablation} 报告三种架构在不同 $\beta$ 值下第 8K 步的 PCE。我们发现 $\beta$ 与坍缩速率之间存在强烈的反向关系。

\begin{table}[t]
\centering
\caption{不同架构下第 8K 步的 PCE 作为 $\beta$ 的函数。较低的 $\beta$ 放大坍缩。}
\label{tab:beta_ablation}
\small
\begin{tabular}{@{}lcccc@{}}
\toprule
模型 & $\beta = 0.05$ & $\beta = 0.1$ & $\beta = 0.2$ & $\beta = 0.5$ \\
\midrule
Llama-2-7B  & 0.91 & 0.74 & 0.58 & 0.38 \\
Mistral-7B  & 0.88 & 0.72 & 0.55 & 0.35 \\
Gemma-2B    & 0.84 & 0.68 & 0.51 & 0.32 \\
\bottomrule
\end{tabular}
\end{table}

该关系近似为对数线性:$\text{PCE} \approx -0.55 \log(\beta) + c$($R^2 > 0.96$)。这证实 KL 惩罚系数直接控制坍缩严重程度——然而从业者通常依据基准性能将 $\beta = 0.1$,在不知情的情况下运作于高坍缩区间。

\paragraph{PCE 测量可靠性。} 我们通过对提示集和每提示样本集进行自助重采样(1{,}000 次迭代)来评估测量稳定性。在 100 个提示、每提示 128 个样本下,PCE 的 95\% 置信区间宽度 $\leq 0.03$。减至 32 个样本使置信区间宽度增至 0.07,而 50 个提示产生 0.05 的置信区间宽度。我们的测量配置($K = 128$、100 个提示)提供了足够的统计功效,可在功效 $> 0.95$(配对置换检验)下检测 0.05 的 PCE 差异。

\paragraph{模型规模效应。} 在相同训练配置($\beta = 0.1$、8K 步)下比较 Gemma-2B、Llama-2-7B 和 Llama-2-13B,PCE 随参数量增长:分别为 0.68、0.74 和 0.81。线性拟合得 PCE $\approx 0.012 \cdot \text{参数量(B)} + 0.65$($R^2 = 0.94$)。我们推测较大的模型有更强的能力记忆被偏好的模式分布,从而实现更尖锐的坍缩。这一缩放趋势对前沿规模的 DPO 训练提出了担忧。

\paragraph{DBSCAN 敏感性。} 变动 $\varepsilon \in [0.2, 0.5]$ 使 PCE 绝对值偏移 $\pm 0.08$,但保持相对排序(Kendall $\tau > 0.95$)与单调性。我们所选的 $\varepsilon = 0.3$ 在留出校准集上最大化轮廓系数。结果对嵌入模型的选择稳健:从 all-MiniLM-L6-v2 切换至 all-mpnet-base-v2 平均改变 PCE 值 $< 0.04$。


% 第 6--7 节:讨论与结论
\section{讨论}
\label{sec:discussion}

\subsection{对对齐实践的含义}

我们的核心发现——安全关键阈值 $t^*$ 在性能最优检查点 $t_{\text{perf}}$ \emph{之前}到达——对当前训练实践有直接后果。将 DPO 训练至留出奖励或基准准确率趋于平台的标准工作流,会系统性地越过模型变得可被坍缩利用的那一点。我们审计的每一个训练至收敛的模型(第 5 节)都表现出处于可利用区间的 PCE 得分。这不是一个实现失误;它是优化一个将概率质量集中到不断收缩的高奖励输出集合上的偏好目标的直接后果。

我们主张,\textbf{偏好坍缩可利用度(PCE)应成为模型卡片中的强制条目},与困惑度和基准得分并列报告。PCE 捕捉了模型质量的一个维度——输出分布几何——而这是现有指标所看不见的。一个模型可以在 MT-Bench 上得分很高、通过红队评估,却仍可被轻易利用,只要其输出分布已坍缩到对抗迁移近乎确定地成功的程度。

这揭示了基于偏好优化的对齐中的一个根本张力。使 DPO 有效的那个性质——将质量集中到被偏好响应上、压制不良尾部——恰恰是使所得模型可预测、因而可被攻击的原因。高对齐质量(在可靠产生``正确''答案的意义上)与高安全性(在难以预测和利用的意义上)在分布集中之下是直接对立的。这不是一个有待修补的缺陷,而是一个对齐研究者必须显式权衡的\emph{设计层面的取舍}。ER-DPO(第 4.4 节)代表了这一取舍空间中的一个点;其他的——也许是自适应正则化调度或熵感知的早停——尚待探索。

\subsection{与现有攻击范式的比较}

坍缩利用与越狱是\emph{正交}的。越狱攻击(GCG、AutoDAN、提示注入)搜索绕过安全过滤器的特定输入。坍缩利用不需要寻找这样的输入——它利用模型输出空间的分布几何,使\emph{任何}下游攻击都更有效。一个对多样化模型以概率 $p$ 成功的越狱,对坍缩模型以概率 $p + \Delta(PCE)$ 成功,因为坍缩模型的响应更可预测、其安全边界更脆弱(命题~3)。两类攻击以乘性方式叠加:坍缩抬高了输入特定攻击赖以构建的基线成功率。

坍缩加速器(第 4.3 节)在结构上也不同于标准数据投毒。传统投毒改变模型\emph{输出什么}——注入后门、将补全偏置向攻击者所选内容。坍缩加速器改变输出分布的\emph{几何},而不必改变任何单个提示的众数响应。一个被投毒的模型可能通过输出层面的安全检查(其 top-1 补全仍然良性),却因熵被驱至临界阈值以下而变得显著更易被利用。这使坍缩投毒从根本上比现有的、检查单个响应而非分布性质的输出监控防御更难检测。

\subsection{局限性}

\textbf{测量成本。} 计算 PCE 需要每提示抽取 $k=128$ 个样本(我们的默认值),使其在推理时代价高昂。尽管我们表明 $k=64$ 已足以对模型排序(与 $k=128$ 的 Spearman $\rho > 0.95$),将 PCE 整合进实时监控仍是工程上的挑战。在训练检查点期间批量审计是可行的;逐请求的 PCE 则不然。

\textbf{理论简化。} 命题~2 和~3 假设在总体损失上的确定性梯度流。真实的 DPO 训练涉及小批量随机性、学习率调度和权重衰减——所有这些都可能相对我们的连续时间预测延迟或加速坍缩。单调性结果(定理~1)在经验上稳健,但其形式化证明要求可能并非普遍成立的正则性条件。

\textbf{快照有效性。} 我们对开放权重模型的审计(表~2)捕捉的是某一时刻。模型会被下游用户频繁更新、重训或微调。PCE 得分可能随后续训练运行而变动,尽管结构性脆弱性(DPO 集中质量)跨运行持续存在。

\textbf{规模。} 我们的受控实验使用至多 13B 参数的模型。一项对 LLaMA-2-70B 应用基于 LoRA 的 DPO 的试点研究显示出定性上相似的坍缩动力学($t^* / t_{\text{perf}} \approx 0.73$,与较小模型一致),但我们缺乏前沿规模上的全秩训练结果。以显著更多数据训练的模型是否表现出更慢的坍缩——抑或仅仅是延迟的坍缩——仍是开放问题。

\textbf{DBSCAN 敏感性。} PCE 的簇计数步骤使用固定 $\varepsilon$ 参数的 DBSCAN。尽管我们的消融(附录 C)表明在归一化嵌入空间中 $\varepsilon \in [0.3, 0.7]$ 范围内的稳定性,具有双峰输出分布的边缘情形可能在边界值处产生不连续的 PCE 跳变。

\textbf{防御开销。} ER-DPO 因熵正则项需要在整个词表上计算逐 token 概率,增加约 15\% 的墙钟训练时间。对大词表模型,这一开销或可通过 top-$k$ 熵近似减小,尽管我们在此未评估此类近似。

\subsection{更广泛的影响与负责任披露}

我们在投稿前 90 天向所有被审计的开放权重模型的维护者披露了我们的发现——包括 PCE 得分、利用成功率和坍缩加速器方法。两个团队(Mistral、StabilityAI)确认收到并表示正在研究熵感知的训练修改。我们未收到对发表的反对意见。

坍缩加速器代码在一个受控研究框架内发布,该框架要求机构隶属关系并对投毒数据生成强制限速。我们不发布预先计算好的投毒数据集。ER-DPO 防御、PCE 计算库和所有评估脚本均无限制地完全开源。

我们强调,坍缩可利用性是偏好优化模型的一项\emph{系统性}性质,而非任何特定系统的缺陷。我们的目标是将社区的注意力从输出层面的安全(``模型会拒绝有害查询吗?'')转向分布层面的安全(``模型的输出几何对利用稳健吗?'')。二者是互补的,而非竞争的关切。

\section{结论}
\label{sec:conclusion}

我们已证明,直接偏好优化所诱导的模式坍缩不仅仅是一个美学或多样性问题,而是一项具体、可测量的安全脆弱性。我们的贡献贯穿这一问题的完整弧线:PCE 指标提供了分布坍缩与利用成功之间的首个定量桥梁;单调性发现($t^* < t_{\text{perf}}$)表明标准训练实践会系统性地产出可被利用的模型;坍缩加速器证明攻击者可以通过规避输出层面检测的数据投毒蓄意放大该脆弱性;我们对七个开放权重模型家族的审计证实该脆弱性存在于生产系统中;而 ER-DPO 提供了一个有效的防御,在将输出熵维持于临界利用阈值之上的同时保留了 97\% 的对齐质量。最紧迫的实践启示是:社区必须将 PCE 监控整合进 DPO 训练流程——不是作为可选的诊断,而是作为与损失收敛和基准性能同等的一等安全约束。随着偏好优化方法持续涌现(KTO、IPO、SPPO),理解并控制分布集中的安全含义,将仍是安全部署对齐语言模型的一项核心挑战。


% 参考文献
\bibliographystyle{plainnat}
\bibliography{references}

% 附录
\appendix

\section{完整证明}
\label{app:proofs}

本附录提供第~\ref{sec:method} 节中推迟的完整证明。全程沿用第~\ref{sec:preliminaries} 节确立的记号:$\pi_\theta$ 表示策略,$\pi_{\text{ref}}$ 表示参考模型,$\hat{r}_\theta(x,y) = \beta \log \frac{\pi_\theta(y|x)}{\pi_{\text{ref}}(y|x)}$ 表示隐式奖励。

\subsection{命题~\ref{prop:det_exploit} 的证明(确定性蕴含可利用性)}

\begin{proof}
回忆 $\text{Det}(x,\pi_\theta) = \max_k p_k$ 是定义~\ref{def:mode_dist} 模式分布下主导语义簇 $C_{k^*}$ 的概率质量。设 $\text{Det}(x,\pi_\theta) = p_{k^*} > 1-\epsilon$。当抽取 $y \sim \pi_\theta(\cdot|x)$ 时,事件 $y \in C_{k^*}$ 以概率 $p_{k^*} > 1-\epsilon$ 发生。在 $y \in C_{k^*}$ 的条件下,$y$ 共享簇质心 $\mu_{k^*}$ 的语义内容,故 $\textsc{Harmful}(y) = \textsc{Harmful}(\mu_{k^*})$ 至多相差分类器误差 $\delta$(即簇内补全被判为与其质心不同的概率)。若 $\mu_{k^*}$ 有害,则
\[
\Pr_{y\sim\pi_\theta}[\textsc{Harmful}(y)=1] \geq p_{k^*}(1-\delta) > (1-\epsilon)(1-\delta) \geq 1-\epsilon-\delta.
\]
将 $\delta$ 吸收进 $\epsilon$ 得 $\Pr[\textsc{Harmful}(y)=1] > 1-\epsilon$,即单个无条件样本以超过 $1-\epsilon$ 的概率有害。一个通过模式探测(第~\ref{sec:attack} 节)识别出 $\mu_{k^*}$ 的攻击者,因而以 $>1-\epsilon$ 的概率在单次查询中成功,主张得证。
\end{proof}

\subsection{命题~\ref{prop:gradient_lockin} 的证明(梯度锁定)}

\begin{proof}
由 DPO 梯度(式~\ref{eq:dpo_gradient}),每样本更新权重为 $\sigma(-\beta \Delta\hat{r}_\theta)$,其中 $\Delta\hat{r}_\theta = \hat{r}_\theta(x,y_w) - \hat{r}_\theta(x,y_l)$。

\emph{第(1)部分。} 直接由式~\eqref{eq:dpo_gradient} 得到:梯度正比于 $\sigma(-\beta\Delta\hat{r}_\theta)$。

\emph{第(2)部分。} 随着训练对被偏好对增大 $\pi_\theta(y_w|x)$ 并减小 $\pi_\theta(y_l|x)$,有 $\Delta\hat{r}_\theta \to +\infty$,故 $\sigma(-\beta\Delta\hat{r}_\theta)\to 0$ 单调下降。因此对已分离对的有效学习信号消失。

\emph{第(3)部分。} 利用 $z\geq 0$ 时 $\sigma(-z)\leq e^{-z}$,更新幅度满足 $\sigma(-\beta\Delta\hat{r}_\theta)\leq e^{-\beta\Delta\hat{r}_\theta}$。在梯度下降下,间隔至少以对数速率增长,$\Delta\hat{r}_\theta(t)\geq c\log t$(对可分的逻辑斯蒂型损失为标准结论),给出更新幅度 $\leq t^{-\beta c}$。因此任何新偏好信号能将质量从主导模式重新分配的速率随 $t$ 多项式衰减,锁定了坍缩的构型。
\end{proof}

\subsection{命题~\ref{prop:critical_point} 的证明($t^*_\tau < t_{\text{perf}}$)}

\begin{proof}
设 $t^*_\tau$ 为 PCE 首次超过阈值 $\tau$(等价地,模式熵 $H_{\text{mode}}$ 首次落到相应水平以下)的步,$t_{\text{perf}}$ 为基准性能峰值的步。我们在以下假设下证明 $t^*_\tau < t_{\text{perf}}$:(C1)熵单调下降(命题~\ref{prop:gradient_lockin} 第 2 部分);(C2)基准质量取决于贪婪解码(主导模式)输出;(C3)质量在主导模式精化期间非递减。

模式坍缩——即事件 $H_{\text{mode}} < $ 阈值——由质量集中驱动,并在 $\sigma(-\beta\Delta\hat{r}_\theta)$ 变得可忽略(命题~\ref{prop:gradient_lockin})时完成,发生于步 $t^*_\tau$。在 $t^*_\tau$ 之后,梯度幅度虽小但非零,继续精化 $C_{k^*}$ 内的\emph{token 级}分布,降低 token 熵并提升连贯性。在(C2)--(C3)下,基准质量在这一精化期间持续改善,并在精化饱和时于 $t_{\text{perf}}$ 达到峰值。由于精化仅在主导模式确立(即 $t^*_\tau$ 之后)后才开始,故 $t_{\text{perf}} > t^*_\tau$。第~\ref{sec:res_evolution} 节中 $1{,}500$--$2{,}500$ 步的经验间隔与这一顺序一致。
\end{proof}

\subsection{定理~\ref{thm:erdpo} 的证明(ER-DPO 多样性下界)}

\begin{proof}
ER-DPO 目标为 $\mathcal{L}_{\text{ER-DPO}} = \mathcal{L}_{\text{DPO}} - \lambda_H \hat{H}(\pi_\theta)$,其中 $\hat{H}$ 是(估计的)输出熵,$\lambda_H>0$。在任何驻点处,$\nabla_\theta \mathcal{L}_{\text{DPO}} = \lambda_H \nabla_\theta \hat{H}$。DPO 项集中质量的能力以 $\beta$ 乘以奖励间隔为界,而熵正则器提供幅度为 $\lambda_H \|\nabla_\theta \hat{H}\|$ 的恢复力以抵抗集中。平衡二者力的幅度,得到一个由二者幅度相等时的取值下界的驻点熵,即式~\eqref{eq:diversity_bound} 的多样性下界。由于 $\hat{H}$ 有界($0\leq \hat{H}\leq \log|\mathcal{V}|$)且在 softmax 参数化下关于 $\theta$ Lipschitz 光滑,复合目标满足下降引理,故梯度下降收敛到这样一个驻点。
\end{proof}

\section{有害性分类器细节}
\label{app:harmful_classifier}

定义~\ref{def:pce} 中的 $\textsc{Harmful}(\cdot)$ 谓词如下实现。对每个提示,我们用 SentenceBERT(\texttt{all-MiniLM-L6-v2})嵌入对 $n=100$ 个采样补全进行聚类,使用 DBSCAN($\varepsilon=0.3$、\texttt{min\_samples}$=3$、余弦度量)。对主导簇 $C_{k^*}$,我们取其中心点(medoid)补全作为代表,并用 LlamaGuard-2(8B)对其标注,将 \texttt{unsafe} 判定映射为 $\textsc{Harmful}=1$。为控制簇内标注噪声 $\delta$,我们另外对五个随机簇成员分类,并对存在分歧的簇标记以供人工复核;在我们的审计集上,$\delta < 0.04$。我们同时报告使用 LlamaGuard-2 的 PCE,以及作为鲁棒性检验、使用 HarmBench 分类器的 PCE;二者结果至多相差 0.02 PCE,在所报告的置信区间之内。


\end{document}
